% Independent book-end source notes; no translation fact or source formula is changed.
\chapter*{勘误与原文疑点}
\phantomsection
\addcontentsline{toc}{chapter}{勘误与原文疑点}
\markboth{勘误与原文疑点}{勘误与原文疑点}

本表记录指定英文版本中的待核验疑点、待补项和省略说明。各条只表示观察到的原文字样，
不表示已判定数学结论错误；译文仍忠实保留该版本的数学内容。

英文版本：\href{https://github.com/stacks/stacks-project/commit/a04446e57ec1fbc252a871afcec7752fb2807b14}{The Stacks Project，2026-07-28}。
\par\noindent
Git commit：\nolinkurl{a04446e57ec1fbc252a871afcec7752fb2807b14}。
永久 Tag 链接便于定位；版本证据以本表列明的锁定 commit 为准。

\begingroup
\small
\setlength{\LTleft}{0pt}
\setlength{\LTright}{0pt}
\renewcommand{\arraystretch}{1.2}
\begin{longtable}{@{}p{0.16\textwidth}p{0.63\textwidth}p{0.13\textwidth}@{}}
\hline
位置 & 原文观察与状态 & 留档 \\
\hline
\endfirsthead
\hline
位置 & 原文观察与状态 & 留档 \\
\hline
\endhead
\hline
\endfoot
\href{https://stacks.math.columbia.edu/tag/0032}{0032}\newline
categories\newline Section 002Z\newline 证明
& \textbf{待核验。}原文先在积范畴（product category）$\mathcal{I}\times\omega$
  中取有限生成子范畴（finitely generated subcategory）$\mathcal{F}$，
  再在同一环境中构造余锥（cocone）及扩充 $\mathcal{F}^{+}$；
  随后称 $\mathcal{F}^{+}$ 为 $\mathcal{I}$ 的有限生成子范畴（finitely generated subcategory）。
  可能缺少投影（projection）说明或环境的记号说明，尚未作数学判定。
  译文保留原文的 $\mathcal{I}$ 表述。
& \href{https://github.com/dongyaotalk/stacks-project-zh/issues/691}{\#691} \\
\hline
\href{https://stacks.math.columbia.edu/tag/04VA}{04VA}\newline
sets\newline Section 000H\newline 引理与证明
& \textbf{原文待补／省略。}引理第三项为
  \texttt{add more here as needed.}，证明中有 \texttt{Details omitted.}。
  本条记录原文待补项和省略细节的事实；译文未添加性质或证明。
& \href{https://github.com/dongyaotalk/stacks-project-zh/issues/698}{\#698} \\
\hline
\href{https://stacks.math.columbia.edu/tag/0AHL}{0AHL}\newline
sets\newline Section 000H\newline 证明
& \textbf{待核验。}原文先定义 $W_{a,i}$，末尾余积（coproduct）公式使用 $W_{i,a}$。
  本条记录下标顺序差异，尚未判定其对论证的影响；译文保留原公式的下标顺序。
& \href{https://github.com/dongyaotalk/stacks-project-zh/issues/699}{\#699} \\
\hline
\end{longtable}
\endgroup

原文证据及完整语义坐标保存在对应留档 Issue 和项目开发文档中。
英文版本变更后须重新核验各条；获得确认的勘误仍需单独记录确认依据及适用版本。
